\documentclass[11pt]{article}
\usepackage{fontspec}
\usepackage[french]{babel}
\usepackage{amsmath}
\usepackage{amssymb}
\usepackage{amsthm}
\usepackage{mathrsfs}
\usepackage{tikz-cd}
\usepackage{enumitem}
\usetikzlibrary{babel}
\newcommand{\printedpage}[1]{\marginpar{\footnotesize p.~#1}}
\newcommand{\msp}[1]{\marginpar{\footnotesize [#1]}}
\newcommand{\Hom}{\operatorname{Hom}}
\newcommand{\Ob}{\operatorname{Ob}}
\newcommand{\bD}{\mathbf{D}}
\newcommand{\bDelta}{\boldsymbol{\Delta}}
\begin{document}

\printedpage{34}
% head: A. Grothendieck : Dérivateurs, Ch. II, Édition des Universités de Montpellier II et Paris VII
Ceci montre l'équivalence de d, d$'$, qui se voit aussi par l'equivalence de chacune de ces deux conditions avec d$''$, par le critère de cofinalité du paragraphe précedent. Soit D cette condition sur $u$. On a par (4.8) D $\Longrightarrow$ A, donc % [sic] l'equivalence -> l'équivalence, accent missing; précedent -> précédent, accent missing
\[
\mathrm{D} \Longrightarrow \mathrm{A} \Longrightarrow \mathrm{c},
\]
il reste à prouver que c $\Longrightarrow$ D, plus précisément c $\Longrightarrow$ d$''$ et c$'$ $\Longrightarrow$ d$''$ \texttt{[c$'$?]}. Explicitons le morphisme d'adjonction (pour $i$ dans $I$)
\[
u_!\underbrace{u^*i} \longrightarrow i,
\]
i.e. pour $j$ dans $I$
\[
\underbrace{(u_!u^*i)(j)}_{\pi_0(j\backslash(I'/u^*(i)))} \longrightarrow i(j) = \Hom(j,i).
\]

\msp{page 54}Or $I'/u^*(i) \simeq I'/i$, donc on a la flèche
\[
(u_!u^*(i))(j) \simeq \pi_0(j\backslash(I'/i)) \longrightarrow \Hom_I(j,i). \tag*{(4.10)}
\]
Or on voit aussitôt que pour $\alpha : j \longrightarrow i$ dans $\Hom(j,i)$, la fibre de (4.10) en $u$ s'identifie à $\pi_0(j\backslash\underset{\alpha}{I'}/i\ )$, de façon plus précise, $j\backslash\underset{\alpha}{I'}/i$\ \ est la sous-catégorie ouverte et fermée de $J = j\backslash(I'/i)$ correspondant à la partie du $\pi_0(J)$ qui est au dessus de $\alpha$ dans (4.10). Cela achève la démonstration de 4.1.

NB Une autre façon de prouver que c $\Longrightarrow$ A c'est de noter que les deux membres dans le morphisme d'adjonction $u_!u^*F \longrightarrow F$ dans $I^{\wedge}$ sont des foncteurs $I^{\wedge} \longrightarrow I^{\wedge}$ qui commutent aux $\varinjlim$ quelconques. Donc pour vérifier que c'est \texttt{[un]} isomorphisme, il suffit de le vérifier pour $F$ dans $I$, puisque $I$ engendre $I^{\wedge}$ par $\varinjlim$ . Par contre, un argument de ce type ne semble pas prouver que pour vérifier que l'on a un \textit{isomorphisme} $F \longrightarrow u_*u^*F$, il suffise de le faire pour $F$ représentable, i.e. de vérifier b$''$ pour $F$ représentable.

Il resterait à voir si c$'$ $\Longrightarrow$ d$''$ \texttt{[c$'$?]}.

\msp{page 55}Explicitons le morphisme d'adjonction
\[
i \longrightarrow u_*u^*i
\]
(pour $i \in \Ob I$) par ses valeurs sur $j$
\[
\underbrace{i(j)}_{\Hom(j,i)} \longrightarrow \Gamma(I'/j, u^*i). \tag*{$(\ast)$}
\]
Soit donc
\[
u_j : I'/j \longrightarrow I,
\]
% foot: Version 8 juillet 2001
\printedpage{35}
% head: A. Grothendieck : Dérivateurs, Ch. II, Édition des Universités de Montpellier II et Paris VII
on a
\[
\begin{aligned}
\Gamma(I'/j, \underbrace{u^*(i)|I'/j}_{u_j^*(i)}) &\simeq \Hom(e_{I'/j}, u^*(i)) \\ % ?? uncertain: extent of the underbrace -- over "u^*(i)|I'/j" or including the closing parenthesis
&\overset{\mathrm{adj}}{\simeq} \Hom(u_{j,!}e_{I'/j}, i),
\end{aligned}
\]
et la flèche $(\ast)$, fonctorielle en $i$, est déduite d'une flèche
\[
u_{j,!}(e_{I'/j}) \longrightarrow j. \tag*{$(\ast\ast)$}
\]
On voit tout de suite, comme ci-dessus, que $(\ast\ast)$ est \texttt{[un]} isomorphisme si et seulement si pour tout $\alpha : k \longrightarrow j$, la catégorie $k\backslash\underset{\alpha}{I'}/j$ est 0-connexe, et si on l'exige pour tout $j$, $k$, $\alpha$, c'est la condition d$''$. Mais il n'est pas clair que pour qu'on ait $(\ast\ast)$ isomorphisme, il suffise que le Hom des deux membres à valeurs dans tout $i$ dans $I$ (et non tout $F$ dans $I^{\wedge}$) soit \texttt{[un]} isomorphisme. La difficulté provient du fait que $u_{j,!}(e_{I'/j})$ n'est pas forcément représentable, et qu'un $F$ quelconque dans $I^{\wedge}$ \msp{page 56}n'est pas forcément $\varprojlim$ de foncteurs représentables. Donc il est douteux que c$'$ \texttt{[c$'$?]} suffise à entraîner que $u$ est $W_0$-connexe.

\textbf{Corollaire 4.2.} \textit{Pour que $u$ soit $W_0$-connexe, il} suffit \textit{qu'il soit une $W_0$-équivalence universelle, i.e. (cf. 3.1) que les fibres soient \textup{0}-connexes, et que pour toute flèche non identique dans $I$, définissant $\bDelta^1 \longrightarrow I$, $I' \times_I \bDelta^1$ soit \textup{0}-connexe.}

4.3. Cette condition de $W_0$-équivalence universelle est cependant \textit{strictement} plus forte que la $W_0$-connexité, comme on le voit sur l'exemple où $u$ est une \textit{équivalence de catégories} non surjective sur les objets (p.ex. l'inclusion d'un sommet dans un groupoïde 1-connexe ayant plus d'un objet). Plus généralement, tout foncteur de localisation est $W_0$-connexe, alors que ses fibres ne sont pas forcément connexes, et a fortiori ce n'est pas une $W_0$-équivalence universelle.

Cependant, si tout endomorphisme $p$ d'un $i$ dans $I$, tel que $p^2 = p$ (projecteur) est l'identité, et si \msp{page 57}$I' \longrightarrow I$ est `transportable', p.ex. si dans $I$ tout isomorphisme est une identité, alors $u$ $W_0$-connexe implique que les fibres de $u$ sont 0-connexes, comme on voit en exprimant d$''$ pour $\alpha$ flèche identique. Les conditions envisagées sur $I$ sont satisfaites si $I$ est ordonnée. Mais même si $I'$ et $I$ sont toutes deux ordonnées, $u$ peut être $W_0$-connexe, sans être une $W_0$-équivalence universelle. Voici un exemple avec $I = \bDelta^2$ :
\[
\begin{array}{r@{\;}l@{\quad}c}
I' = & \left\{\begin{tikzcd}[ampersand replacement=\&]
 \& 1' \arrow[r] \arrow[d] \& 2 \\
0 \arrow[r] \& 1 \&
\end{tikzcd}\right\} & \\[1ex]
 & & \Big\downarrow \\[1ex]
I = & \left\{\begin{tikzcd}[ampersand replacement=\&]
0 \arrow[r] \& 1 \arrow[r] \& 2
\end{tikzcd}\right\}. &
\end{array}
\]
% foot: Version 8 juillet 2001

\printedpage{36}
% head: A. Grothendieck : Dérivateurs, Ch. II, Édition des Universités de Montpellier II et Paris VII
Les catégories bilocalisées de $I$ sont les trois sommets $\{0\}$, $\{1\}$, $\{2\}$, et les trois segments $[0,1]$, $[1,2]$ et $[0,2] = I$, les images inverses sont respectivement
\[
\bullet\ ,\quad
\begin{tikzcd}[row sep=normal] {} \arrow[d] \\ {} \end{tikzcd}\ ,\quad
\bullet\ ,\quad
\begin{tikzcd}[row sep=normal] & {} \arrow[d] \\ {} \arrow[r] & {} \end{tikzcd}\ ,\quad
\begin{tikzcd}[row sep=normal] {} \arrow[r] \arrow[d] & {} \\ {} & \end{tikzcd}\ ,\quad
I'\quad ,
\]
elles sont toutes 0-connexes, i.e. $W_0$-asphériques, et même contractiles donc $W_\omega$-asphériques (donc $u$ est même $W_\omega$-connexe, i.e. $W_\omega$-biasphérique). Mais ce n'est pas une $W_0$-équivalence universelle, car l'image inverse de $I'$ par $\bDelta^1 \xrightarrow{\partial} \bDelta^2$, $\partial(0) = 0$, $\partial(1) = 2$, \msp{page 58}est la catégorie disconnexe $\varepsilon = \{\bullet\ \ \bullet\}$.

J'avais espéré un moment que les foncteurs $W_0$-connexes $I' \xrightarrow{u} I$ seraient ceux qui se factorisent en $I' \xrightarrow{u_0} I'' \xrightarrow{v} I$, avec $u_0$ une équivalence, et $v$ une $W_0$-équivalence universelle. Mais l'exemple précédent le met en défaut, avec $u$ même $W_\omega$-connexe. En effet, $I''$ serait déduit d'un ensemble préordonné, et $v$ se factoriserait par l'ensemble ordonné correspondant, savoir $I'$, donc serait de la forme
\[
I'' \longrightarrow I' \xrightarrow{u} I,
\]
où $I'' \longrightarrow I'$ est un passage au quotient, et l'image inverse de $\bDelta^1 \xrightarrow{\partial} I$ serait encore disconnexe (un $\pi_0$ de cardinal 2). Donc c'est vraiment sans espoir.

Cependant

\textbf{Corollaire 4.3.} \textit{Supposons que $u$ soit $W_0$-lisse, ou $W_0$-propre, ou une $W_0$-fibration faible. Alors $u$ est $W_0$-connexe si et seulement si c'est une $W_0$-équivalence universelle, i.e. si et seulement si a) ses fibres sont \textup{0}-connexes et b) $I' \times_I \bDelta^1$} \texttt{[est]} \textit{\textup{0}-connexe pour tout $\bDelta^1 \longrightarrow I$ (correspondant à une flèche non identique).}

\msp{page 59}Dans les trois cas on sait que $u$ est une $W_0$-équivalence universelle si et seulement si ses fibres sont $W_0$-connexes. Or la $W_0$-lissité implique que $\forall i \in \Ob I$, l'inclusion
\[
I'_i \longrightarrow i\backslash I'
\]
est une $W_0$-équivalence, d'autre part $i\backslash I' \longrightarrow i\backslash I$ étant $W_0$-asphérique, donc une $W_0$-équivalence par hypothèse, $i\backslash I'$ est 0-connexe, donc $I'_i$ aussi. Le cas $u$ $W_0$-propre est dual. Dans le cas \texttt{[où]} $u$ \texttt{[est]} une $W_0$-fibration faible, on sait encore que $I'_i \longrightarrow i\backslash I'$ est une $W_0$-équivalence, et on conclut comme dans le cas lisse (ou, au choix, via $I'_i \longrightarrow I'/i$).

Mais l'idée me vient que les foncteurs $W_0$-connexes pourraient être exactement les foncteurs de localisation ? Il revient au même de dire qu'un foncteur $u : I' \longrightarrow I$ \textit{conservatif} (i.e. $u(\alpha)$ isomorphisme $\Longrightarrow$ $\alpha$ isomorphisme) qui est $W_0$-connexe, est une équivalence. Mais c'est faux, comme on voit en prenant $I = e$, $I'$ un groupoïde 0-connexe et non 1-connexe : $I' \longrightarrow e$ est $W_0$-connexe (et même une $W_0$-équivalence universelle), mais ce n'est pas une équivalence (et pas un foncteur de localisation).
% foot: Version 8 juillet 2001

\printedpage{37}
% head: A. Grothendieck : Dérivateurs, Ch. II, Édition des Universités de Montpellier II et Paris VII
\msp{page 60}Je voudrais généraliser la notion de flèche $W_0$-biasphérique au cas d'un dérivateur $\bD$ quelconque.

\textbf{Proposition 4.4.} \textit{Soit $\bD$ un dérivateur sur} Diag, \textit{et $u : I' \longrightarrow I$ une flèche dans} Diag. \textit{Supposons que $u_*$ et $u_!$ existent, et se calculent de façon standard, pour tout flèche $u$ dans} Diag. \textit{(En fait, il suffit que ce soit le cas pour $u$ et ses localisées $u/i : I'/i \longrightarrow I/i$, ou $u$ et ses colocalisées $i\backslash u : i\backslash I' \longrightarrow i\backslash I$.) Considérons les conditions :} % [sic] pour tout flèche -> pour toute flèche, gender agreement
\begin{itemize}[leftmargin=4em]
\item[\textit{a)}] $u^* : \bD(I) \longrightarrow \bD(I')$ \texttt{[est]} \textit{pleinement fidèle.}
\item[\textit{a$'$)}] \textit{$\forall$ $F$ dans $\bD(I)$, la flèche}
\[
u_!u^*F \longrightarrow F
\]
\textit{est un isomorphisme.}
\item[\textit{a$''$)}] \textit{$\forall$ $F$ dans $\bD(I)$, la flèche}
\[
u_*u^*F \longleftarrow F
\]
\textit{est un isomorphisme.}
\item[\textit{b)}] \textit{$\forall$ $i$ dans $I$, la flèche $i\backslash I' \longrightarrow i\backslash I$ est $\bD$-asphérique.}
\item[\textit{b$'$)}] \textit{$\forall$ $i$ dans $I$, la flèche $I'/i \longrightarrow I/i$ est $\bD$-coasphérique.}
\item[\textit{b$''$)}] \textit{$\forall$ $\alpha : i_0 \longrightarrow i_1$ dans $I$, la catégorie $i_0\backslash\underset{\alpha}{I'}/i_1$ est $\bD$-asphérique.}
\end{itemize}
\textit{On a alors les implications}
\[
(b \Longleftrightarrow b' \Longleftrightarrow b'') \Longrightarrow (a \Longleftrightarrow a' \Longleftrightarrow a'').
\]
\textit{On a $a \Longrightarrow b$ (i.e. toutes les conditions licites envisagées sont équivalentes)} \msp{page 61}\textit{dans chacun des cas suivants.}
\begin{itemize}[leftmargin=4em]
\item[$1^o$] \textit{$I$ est préordonné.}
\item[$2^o$] \textit{$\forall$ $i$ dans $I$, $\alpha_i^* : \bD(I) \longrightarrow \bD(I/i)$ est essentiellement surjectif.}
\item[$3^o$] \textit{$\forall$ $i$ dans $I$, considérant $\alpha_i : I/i \longrightarrow I$, le foncteur $\alpha_{i,!} : \bD(I/i) \longrightarrow \bD(I)$ est conservatif.}
\item[$4^o$] $\mathcal{D} = \mathrm{HOT}_{W_0} \approx \bD_{\mathrm{Ens}}$.
\end{itemize}
\textsc{Démonstration}. L'équivalence de a, a$'$, a$''$ est claire, soit $A$. De même l'equivalence de b, b$'$, b$''$, soit $B$. % [sic] l'equivalence -> l'équivalence, accent missing

Prouvons
\[
B \Longrightarrow A
\]
en supposant p.ex. l'existence de $u_*$. On admet tacitement, avec l'existence de $u_*$ : $\bD(I') \longrightarrow \bD(I)$, le mode de calcul des fibres de $u_*(-)$. Donc $A$, via a$''$, équivaut à
\[
H^{\bullet}_{\bD}(I'/i, u^*F) \xleftarrow{\mathit{iso}} H^{\bullet}_{\bD}(I/i, F) \simeq F(i)
\]
% foot: Version 8 juillet 2001
\printedpage{38}
% head: A. Grothendieck : Dérivateurs, Ch. II, Édition des Universités de Montpellier II et Paris VII
pour tout $F$ dans $\bD(I)$, tout $i$ dans $I$.

Soit
\[
\alpha_i : I/i \longrightarrow F, \qquad F_i = \alpha_i^*(F), % [sic] I/i \longrightarrow F -> I/i \longrightarrow I, alpha_i is defined as I/i -> I in 3^o
\]
on doit donc exprimer que pour tout $F_i$ provenant d'un $F$, on a
\[
H^{\bullet}_{\bD}(I/i, F_i) \xrightarrow{\mathrm{iso}} H^{\bullet}_{\bD}(I'/i, (u/i)^*(F_i)). \tag*{$(\ast)$}
\]
Or sous réserve d'existence de $(u/i)_!$, dire que $(\ast)$ est \texttt{[un]} isomorphisme pour \textit{tout} $F_i$ de $\bD(I/i)$, signifie que
\[
u/i : I'/i \longrightarrow I/i
\]
\msp{page 62}est $\bD$-coasphérique, d'où b$'$ $\Longrightarrow$ a$''$. Il est clair que l'implication inverse est valable si $\alpha_i^*$ est essentiellement surjectif, c'est le cas $2^o$. Or $2^o$ est satisfait si $I$ est préordonné, d'où le cas $1^o$. Le cas $4^o$ a déjà été traité (4.1). Il faut encore traiter le cas $3^o$, qui contient d'ailleurs le cas $2^o$, donc le cas $1^o$ (et aussi le cas $4^o$, cf. plus bas).

Il faut exprimer l'isomorphie de la flèche $(\ast)$ pour $F_i$ de la forme $\alpha_i^*(F)$. Soit $\xi$ \texttt{[un]} objet de $\mathcal{A} = \mathcal{D}(e)$, il faut exprimer l'isomorphie de $\Hom(\xi, -)$ dans les deux membres de $(\ast)$, i.e. de
\[
\begin{aligned}
\Hom(I/i, F_i) \quad &\longrightarrow \quad \Hom(\xi'_{I'/i}, (u/i)^*(F_i)) \\
&\overset{\text{adj.}}{\simeq} \quad \Hom((u/i)_!(\xi'_{I'/i}), F_i),
\end{aligned}
\tag*{$(\ast\ast)$}
\]
\texttt{[plutôt} $\xi_{I'/i}$\texttt{]} flèche qui se déduit de
\[
(u/i)_!(u/i)^*(\xi_{I'/i}) \xrightarrow{\text{adj.}} \xi_{I'/i} \tag*{$(\ast\ast\ast)$}
\]
en prenant le Hom dans $F_i$. La condition $u/i$ coasphérique (i.e. b$'$) signifie que $(\ast\ast\ast)$ est un isomorphisme pour tout $\xi$ dans $\mathcal{A}$. \textit{Suffit}-il pour ceci que son Hom dans un $F_i$ soit bijectif, quand on restreint $F_i$ dans $\bD(I/i)$ à être dans l'image de $\bD(I)$ par $\alpha_i^*$ ? Or cette flèche $(\ast\ast)$, pour $F_i = \alpha_i^*F$, s'explicite par la formule d'adjonction relative \msp{page 63}au foncteur $\alpha_i$, comme
\[
\Hom(\alpha_{i,!}(\xi_{I/i}), F) \longrightarrow \Hom(\alpha_{i,!}(\ldots), F).
\]
Dire que c'est un isomorphisme pour tout $F$, signifie donc \textit{exactement} que la transformée de la flèche $(\ast\ast\ast)$ par $\alpha_!$ est un isomorphisme. Pour en conclure que $(\ast\ast\ast)$ elle-même est \texttt{[un]} isomorphisme, la condition naturelle à imposer est donc que $\alpha_{i,!} : \bD(I/i) \longrightarrow \bD(I)$ soit conservatif. C'est le cas $3^o$. Cela achève la démonstration de 4.4.

4.5. Montrons que la condition $3^o$ sur $I$ n'est pas toujours satisfaite, dans le cas de $\bD = (\mathrm{HOT}_{W_0})^{\mathrm{opp}} \approx (\bD_{\mathrm{Ens}})^{\mathrm{opp}}$, bien que dans ce cas, il soit pourtant vrai que pour toute $I'$ sur
% foot: Version 8 juillet 2001

\end{document}
